%%\section{\minerva Experiment}

The \minerva detector~\cite{Aliaga:2013uqz} is comprised of a fine-grained scintillator 
central tracking region surrounded by electromagnetic and hadronic calorimeters. 
The core consists of tracking planes made of interleaved scintillator 
strips of triangular profile, enabling charged-particle energy depositions 
to be located to within 3~mm. The planes are mounted vertically, nearly perpendicular 
to the neutrino beam axis which is 58~mrad from horizontal.   
Three different plane orientations ($0^{\circ}$ and $\pm 60^{\circ}$ from the vertical) 
permit three-dimensional reconstruction of charged particle trajectories. 
The detector's 3~ns hit-time resolution allows separation of multiple neutrino interactions 
within each 10~$\mu$s spill from the accelerator. The scintillator planes
are supported by exterior hexagonal steel frames with rectangular 
scintillator bars embedded into slots, which serve as the side hadronic calorimeter. 
The magnetized \minos near detector~\cite{Michael:2008bc} located two meters 
downstream of \minerva serves as a muon spectrometer. 

These data were taken in the \numi beamline at Fermilab when its focusing elements 
were configured to produce an intense beam of muon neutrinos (\numu) peaked at 3.5~GeV. 
The run period for these data occurred between March 2010 and April 2012, and 
corresponds to $3.04\times{10}^{20}$ protons on target (POT). The neutrino flux is over 
$95\%$ \numu in the peak, with the remainder consisting of 
\numubar, $\nu_e$, and $\bar\nu_e$, and is predicted using a \geant-based model 
constrained by hadron production data~\cite{Alt:2006fr} 
as described in Ref.~\cite{Eberly:2014mra}.  

Neutrino interactions are simulated using the \genie 2.6.2~\cite{Andreopoulos:2009rq} 
event generator. The propagation of particles in the detector and the corresponding detector response  
are simulated with \geant~\cite{Agostinelli:2002hh}. The calorimetric 
energy scale is tuned using through-going muons to ensure that the photon
statistics and reconstructed energy deposition agree between simulation and data.
Measurements made with a smaller version of the \minerva detector in a low energy 
hadron test beam~\cite{Aliaga:2013uqz} are used to constrain the 
uncertainties associated with the detector response to both protons and charged pions.  


